\documentclass[./pmain.tex]{subfiles}
\begin{document}



\portadaenet{Diagrama Causa-Efecto}{Mantenimiento y Reparación de Equipos}

\imgfr{presentaciones/images/p2/pezcado}{Diagrama de Ishikawa }
\textfr{Diagrama causa efecto}{
El problema analizado puede provenir de diversos ámbitos como la salud, calidad de productos y servicios, fenómenos sociales, históricos, organización, etc. A este eje horizontal van llegando líneas oblicuas –como las espinas de un pez– que representan las causas valoradas como tales por las personas participantes en el análisis del problema. A su vez, cada una de estas líneas que representa una posible causa, recibe otras líneas paralelas a la línea central que representan las causas secundarias. Cada grupo formado por una posible causa primaria y las causas secundarias que se le relacionan forman un grupo de causas con naturaleza común. Este tipo de herramienta permite un análisis participativo mediante grupos de mejora o grupos de análisis, que mediante técnicas como por ejemplo la lluvia de ideas, sesiones de creatividad, y otras, facilita un resultado óptimo en el entendimiento de las causas que originan un problema, con lo que puede ser posible la solución del mismo.
}

\itfr{Creación del diagrama de causa y efecto}{presentaciones/images/p2/pezcado}{
El procedimiento consta de cinco pasos consecutivos.
\begin{itemize}
\item Dibuje el diagrama de Ishikawa e ingrese los principales factores de influencia
\item  Calcula las causas principales y secundarias
\item  Verifique que esté completo
\item  Selección de declaraciones probables
\item  Comprobando la causa más probable de la corrección
con una prueba de significancia (prueba de hipótesis).
\end{itemize}
}
\textfr{Dibuje el diagrama de Ishikawa e ingrese los principales factores de influencia}{
El punto de partida es una flecha horizontal que apunta hacia la derecha, en cuya punta se encuentra el objetivo o problema que se formula de la manera más sucinta posible. Las flechas de las principales variables de influencia que conducen a un efecto determinado se encuentran en ángulo.
}
\itfr{Calcula las causas principales y secundarias}{presentaciones/images/p2/arbol}{
Las causas potenciales se exploran utilizando técnicas de creatividad. Estos se muestran en forma de flechas más pequeñas en la línea de las respectivas variables de influencia principales. Si existen otras causas subyacentes, puede diversificarse más; esto da como resultado ramificaciones cada vez más finas.
}
\itfr{  Verifique que esté completo}{presentaciones/images/p2/problema}{
Compruebe que se hayan considerado realmente todas las causas posibles. La visualización a menudo facilita la búsqueda de otras causas.
}

\itfr{ Selección de declaraciones probables}{presentaciones/images/p2/importancia}{
Las causas potenciales se ponderan según su importancia e influencia en el problema. También se determina la causa con mayor probabilidad.
}

\itfr{Comprobando la causa más probable de la corrección}{presentaciones/images/p2/causa}{
Sobre la base del conocimiento y la experiencia de los especialistas, finalmente se analiza si realmente se ha determinado la causa correcta del problema. Estadísticamente, la suposición de que la causa identificada es una causa principal se puede respaldar con una prueba de significancia (prueba de hipótesis).
}
\textfr{Factores de influencia}{
El diagrama de Ishikawa se construye con mayor frecuencia utilizando los 5 factores, grupos principales de factores de influencia.
\begin{itemize}
\item El ambiente. Las condiciones ambientales que influyen en el proceso, como el tiempo, la temperatura, la humedad o la limpieza.
\item Factor humano. Cualquier persona que participe en el proceso a lo largo de toda la cadena de producción, incluidas todas las funciones de apoyo.
\item Método. Define cómo se lleva a cabo el proceso y qué requisitos se necesitan para ello, como procedimientos de calidad, órdenes de trabajo, instrucciones de trabajo o planos.
\item Maquinaria. Todas las máquinas y equipos necesarios para realizar la tarea, incluidas las herramientas.
\item Materiales. Las materias primas, las piezas compradas y los conjuntos que entran en el producto final.

\end{itemize}
}
\itfr{El ambiente}{presentaciones/images/p2/ambiente}{Las condiciones ambientales que influyen en el proceso, como el tiempo, la temperatura, la humedad o la limpieza.}

\itfr{Factor humano.}{presentaciones/images/p2/causa}{Cualquier persona que participe en el proceso a lo largo de toda la cadena de producción, incluidas todas las funciones de apoyo.}

 \itfr{Método.}{presentaciones/images/p2/metodo}{Define cómo se lleva a cabo el proceso y qué requisitos se necesitan para ello, como procedimientos de calidad, órdenes de trabajo, instrucciones de trabajo o planos.}
\itfr{ Maquinaria.}{presentaciones/images/p2/maquinaria}{ Todas las máquinas y equipos necesarios para realizar la tarea, incluidas las herramientas.}
\itfr{ Materiales.}{presentaciones/images/p2/materiales}{ Las materias primas, las piezas compradas y los conjuntos que entran en el producto final.
}



\end{document}